\documentclass{article}
\usepackage{amsmath}
\usepackage{amsfonts}

\begin{document}

\section{Proof of Blurring in Soft Stacks}

\textbf{Goal:} To prove that as stack operations are repeatedly applied, the distribution of elements stored in the stack becomes increasingly high-entropy, thereby losing deterministic fidelity.

\subsection*{Definitions and Notation}
Let $t \in \mathbb{N}$ represent the current discrete time step, and $K \in \mathbb{N}$ denote the finite maximum depth of the stack. The total number of distinct stack symbols is given by $\Gamma \in \mathbb{N}$, where the specific symbol $\Gamma$ designates the empty stack symbol, which cannot be pushed. We use $e_j$ to represent a standard one-hot basis vector; for instance, $e_\gamma$ is the one-hot embedding of stack symbol $\gamma$, and $e_\Gamma$ represents the empty stack symbol. The standard probability simplex over $N$ elements is denoted by $\Delta([N])$. The controller's output vector representing the probability distribution over memory actions at time $t$ is $a(t) \in \Delta([\Gamma+1])$. From this, $a_{push}(t) = \sum_{\gamma=1}^{\Gamma-1} a_\gamma(t)$ is the total probability of pushing any non-empty symbol, $a_{pop}(t) = a_\Gamma(t)$ is the probability of a pop operation, and $a_{noop}(t) = a_{\Gamma+1}(t)$ is the probability of a no-operation. Finally, the matrix representing the entire stack at time $t$ is $M(t) \in \mathbb{R}^{K \times \Gamma}$, with its $k$-th row, $m_k(t) \in \Delta([\Gamma])$, representing the probability mass function of the symbol stored at the $k$-th stack position.

\subsection{Stack elements and initial state}
When an element is pushed onto the stack, it enters as a probability distribution over the non-empty symbols: $s(t) = \sum_{\gamma=1}^{\Gamma-1} a_\gamma(t)e_\gamma$. Let us track the distribution of a specific element residing at depth $k$, denoted by $m_k(t)$.

\subsection{Sequence of stack operations}
During each time step, the stack's state is recursively updated using the affine linear dynamical system $M(t+1) = A(t)M(t) + B(t)$. The transition matrix is defined as $A(t) = a_{noop}(t)I + a_{push}(t)D + a_{pop}(t)U$, where $I$ is the $K \times K$ identity matrix, and $D, U \in \{0,1\}^{K \times K}$ are the downshift and upshift nilpotent operators, respectively. The term $B(t) = e_1s(t)^T + a_{pop}(t)e_K e_\Gamma^T$ represents the boundary injection terms for pushing a new distribution to the top ($e_1$) and pulling the empty symbol from the bottom ($e_K$).

Unrolling this recursion over a sequence of operations yields:
\begin{equation*}
M(t) = \left(\prod_{s=1}^{t-1}A(s)\right)M(1) + \sum_{s=1}^{t-1}\left(\prod_{u=s+1}^{t-1}A(u)\right)B(s)
\end{equation*}
This demonstrates that the state of our tracked element at any future time is governed by the sequential product of the transition matrices $A(u)$.

\subsection{Blurring}
We can observe how the distribution of our element blurs by examining the row-wise update rule for the interior of the stack, where $k \in [2, K-1]$:
\begin{equation*}
m_k(t+1) = a_{push}(t)m_{k-1}(t) + a_{noop}(t)m_k(t) + a_{pop}(t)m_{k+1}(t)
\end{equation*}

As the action weights $a_{push}(t)$, $a_{noop}(t)$, and $a_{pop}(t)$ are derived from a softmax distribution, they effectively apply fractional shifts. Applying this convex combination recursively causes a discrete-time diffusion. The distinct probability mass of the original element $m_k(t)$ becomes blended with the probability distributions from the adjacent slots $m_{k-1}(t)$ and $m_{k+1}(t)$. Furthermore, as operations continue, pop actions inevitably draw the empty stack symbol $e_\Gamma$ up from the truncated bottom boundary of the stack ($k=K$). 

Therefore, if the element remains in the soft stack over an extended amount of time, its initial probability distribution will continuously average with neighboring pushes and the empty symbol $e_\Gamma$. This diffusion monotonically increases the entropy of the probability mass function $m_k(t)$, transitioning it away from a deterministic (one-hot) state toward a high-entropy mixture. 


\end{document}